\chapter{Runny Nose (Rhinorrhea)}

\section*{Definition}
Excessive discharge of nasal mucus from the nostrils, varying in consistency from clear and watery to thick and purulent, reflecting underlying causes.

\section*{What Happens in Your Body}
Increased nasal secretion from goblet cells and submucosal glands triggered by viral infection, allergic response (IgE-mediated mast cell degranulation), vasomotor instability, or irritation. Clear rhinorrhea suggests allergy or viral causes; purulent discharge suggests bacterial sinusitis.

\section*{Causes}
\begin{itemize}
  \item Common cold (viral URI)
  \item Allergic rhinitis (hay fever)
  \item Sinusitis (acute or chronic)
  \item Vasomotor rhinitis
  \item Cold weather or temperature changes
  \item Spicy foods (gustatory rhinitis)
  \item Nasal polyps
  \item Cerebrospinal fluid leak (clear watery unilateral)
  \item Pregnancy rhinitis
  \item Medication-induced (rhinitis medicamentosa from decongestant overuse)
\end{itemize}

\section*{When to See a Doctor}
See your doctor if:
\begin{itemize}
  \item Symptoms are severe or getting worse over time
  \item Persistent symptoms beyond 10 days, thick yellow/green discharge with facial pain, fever, bloody discharge, or unilateral clear drainage after head trauma
  \item You have other concerning symptoms such as fever, weight loss, or bleeding
\end{itemize}

\section*{Related Symptoms}
\begin{itemize}
  \item Sneezing
  \item Nasal congestion
  \item Sore throat
  \item Cough
  \item Itchy eyes
\end{itemize}
\textbf{Important:} If this symptom persists, your doctor will check for serious underlying causes.

\section*{Self-Care / Home Management}
\begin{itemize}
  \item Rest and avoid activities that make it worse.
  \item Maintain adequate hydration and a balanced diet.
  \item Over-the-counter remedies may provide symptomatic relief (consult a pharmacist or doctor).
  \item Monitor symptoms and seek professional advice if they persist or worsen.
  \item Avoid self-medicating with prescription medications without medical guidance.
\end{itemize}

\section*{How Your Doctor Will Evaluate You}
\begin{itemize}
  \item A thorough history and physical examination are the first steps in evaluation.
  \item Laboratory tests (blood work, urinalysis) may be ordered based on clinical suspicion.
  \item Imaging studies (X-ray, ultrasound, CT, MRI) may be indicated when structural pathology is suspected.
  \item Specialist referral is arranged when the diagnosis remains uncertain or requires specific expertise.
\end{itemize}

\section*{Treatment Options}
\begin{itemize}
  \item Treatment targets the underlying cause identified through diagnostic evaluation.
  \item Supportive care and symptom management are provided while awaiting definitive therapy.
  \item Medications, lifestyle modifications, and physical therapies may be prescribed as appropriate.
  \item Follow-up ensures treatment effectiveness and allows timely adjustments to the care plan.
\end{itemize}

\clearpage